% Modelo anonimizado para submissão a congresso
% Baseado no documento fornecido pelo autor.
\documentclass[12pt,a4paper]{article}

\usepackage[utf8]{inputenc}
\usepackage[T1]{fontenc}
\usepackage[brazil]{babel}
\usepackage{geometry}
\usepackage{setspace}
\usepackage{times}
\usepackage{enumitem}

\geometry{
  top=2.5cm,
  bottom=2.5cm,
  left=2.5cm,
  right=2.5cm
}

\setlength{\parindent}{0pt}
\setlength{\parskip}{6pt}
\onehalfspacing

\begin{document}

\begin{center}

{\large\bfseries IMPACTO DA INTERAÇÃO COM JOGOS DIGITAIS SOBRE O ENGAJAMENTO
COGNITIVO E A ACEITAÇÃO TECNOLÓGICA EM PESSOAS IDOSAS: UMA
ABORDAGEM DE GERONTECNOLOGIA APLICADA À REABILITAÇÃO}

\vspace{0.5cm}

{\normalsize
\textit{Impact of interaction with digital games on cognitive engagement and technological
acceptance in older adults: A gerontechnological approach applied to rehabilitation.}
}

\vspace{0.5cm}

Autor A$^{1}$, Autor B$^{1}$, Autor C$^{1}$,
Autor D$^{2}$, Autor E$^{1}$, Autor F$^{1}$

\vspace{0.2cm}

$^{1}$ Escola de Artes, Ciências e Humanidades da USP, São Paulo, Brasil\\
$^{2}$ Universidade Federal de São Paulo, São Paulo, Brasil

\vspace{0.2cm}

E-mail correspondente: [REMOVIDO]

\end{center}

\textbf{Introdução} A experiência de jogo mostrou-se bem aceita e promoveu engajamento
cognitivo em idosos, independentemente da função cognitiva ou idade. Os resultados
reforçam o potencial dos jogos digitais como recurso complementar na reabilitação e
estimulação cognitiva, especialmente quando desenvolvidos com acessibilidade,
adaptação progressiva e usabilidade adequada, favorecendo maior adesão e participação
durante as atividades terapêuticas.

\textbf{Objetivo} Investigar o engajamento cognitivo de pessoas idosas durante a interação
com jogos digitais, relacionando a experiência subjetiva de jogo com o desempenho
cognitivo global.

\textbf{Método} Estudo transversal realizado com 73 idosos de 60 a 90 anos
(média 72,5 $\pm$ 5,4), sendo 64 mulheres e 9 homens. Foram aplicados o Montreal
Cognitive Assessment (MoCA) para avaliação cognitiva e o Game Experience
Questionnaire (GEQ) para mensurar imersão, afeto, desafio, fluxo e envolvimento
emocional. As comparações entre sexos e faixas etárias (60--74 e 75--90 anos)
utilizaram o teste de Mann-Whitney. As correlações entre cognição e experiência de
jogo foram verificadas pelo coeficiente de Spearman.

\textbf{Resultados} Os participantes demonstraram níveis positivos de engajamento
durante a interação, com predominância de afeto positivo, sensação de competência e
fluxo moderado. Não houve diferenças significativas entre sexos ($p>0,12$) ou faixas
etárias ($p>0,40$). As correlações entre MoCA e dimensões do GEQ foram fracas,
sugerindo que o engajamento no jogo é relativamente independente do nível de
desempenho cognitivo. Idosos com menor pontuação no MoCA mantiveram boa capacidade
de interação e participação ativa.

\textbf{Conclusão} A experiência de jogo foi bem aceita e promoveu engajamento
cognitivo em idosos, independentemente da função cognitiva ou idade. Os achados
reforçam o potencial dos jogos digitais como recurso complementar na reabilitação e
estimulação cognitiva, especialmente quando acessíveis e adaptáveis.

\textbf{Palavras-chave:} Geriatria; jogos eletrônicos de movimento; envelhecimento;
reabilitação; satisfação do paciente.

\end{document}
